\section{Pointwise curvature and its weighted integral}\label{sec:curvature_obstruction}
The exact odd-population constant defines the factor
\[
 \rho_N=\frac{N-1}{N}L_N=\frac{1}{N\operatorname{artanh}(1/N)}<1.
\]
The corresponding pointwise inequality takes the form
\begin{equation}\label{eq:falseoddcurve}
 \psi_{N,k,m}''(t)\le\rho_N\frac{m(N-m)}{4(N-1)}.
\end{equation}
The next proposition gives an exact tilted variance above the right side of \eqref{eq:falseoddcurve}.

\begin{proposition}[Exact exceedance of the constant pointwise bound]\label{prop:obstruction}
For $N=15$, $k=7$, $m=3$, and $t=\log(7/6)$,
\begin{equation}\label{eq:obs_exact}
 \frac{4(N-1)}{m(N-m)}\psi_{N,k,m}''(t)
 =\frac{307887776}{308178025}
 >\frac{675}{676}>\rho_{15}.
\end{equation}
This nearest-balanced population violates \eqref{eq:falseoddcurve} at the stated tilt.
\end{proposition}

\begin{proof}
With $z=7/6$, the tilted hypergeometric weights for $j=0,1,2,3$ are proportional to
\[
 w_j=\binom7j\binom8{3-j}z^j.
\]
They are rational, so both the mean and the variance are rational. Direct summation gives
\[
 \E_tH=\frac{26313}{17555},\qquad
 \operatorname{Var}_t(H)=\frac{197927856}{308178025}.
\]
Since $m(N-m)/[4(N-1)]=9/14$, the first equality in \eqref{eq:obs_exact} follows by multiplying the variance by $14/9$. The first strict inequality follows from
\[
 \frac{307887776}{308178025}-\frac{675}{676}
 =\frac{111969701}{208328344900}>0.
\]
Finally, the power series for $\operatorname{artanh}$ at $x=1/15$ yields
\[
 \operatorname{artanh}(1/15)>\frac1{15}+\frac1{3\cdot15^3},
 \quad\text{hence}\quad
 \rho_{15}<\left(1+\frac1{3\cdot15^2}\right)^{-1}=\frac{675}{676}.
\]
Combining these two inequalities proves \eqref{eq:obs_exact}.
\end{proof}

\begin{proposition}[Integrated curvature identity]\label{prop:integral}
For every centered random variable with finite moment generating function, writing $\psi(t)=\log\E e^{tX}$ gives
\begin{equation}\label{eq:curv_integral}
 \frac{2\psi(t)}{t^2}=2\int_0^1(1-u)\psi''(ut)\,du,
 \qquad t\ne0.
\end{equation}
The proxy objective is therefore a weighted average of curvature along the segment from zero to the chosen tilt.
\end{proposition}
\begin{proof}
Twice integrate $\psi''$ using $\psi(0)=\psi'(0)=0$ and substitute $s=ut$; the formula is valid for either sign of $t$.
\end{proof}

Theorem \ref{odd:thm:envelope} bounds the weighted integral \eqref{eq:curv_integral} through the mean-dependent envelope of Lemma \ref{odd:lem:barrier}. Proposition \ref{prop:obstruction} identifies a point of this variance curve above the constant threshold in \eqref{eq:falseoddcurve}.
